\chapter{Packages and Profiles}
\label{ch:packages}

\standfirst{Code you intend to keep lives in files, and which files compose an
image is one line in a profile. There is no package manager, and this chapter
is why you do not need one.}

\section{The layout}

\begin{plain}
sources/     the 1983 class library, MIT, vendored.  NEVER EDITED.
lib/         ours: exceptions, concurrency, SUnit, protocol shims, SQL, JSON
pharo/       imported Pharo packages, each with a PROVENANCE.md
profiles/    which packages compose an image
\end{plain}

The rule that makes this work: \textbf{\ct{sources/} is never edited}. Every
divergence from 1983 is a new file in \ct{lib/} or \ct{pharo/}. That buys
four things---the Blue Book path stays untouched code rather than
carefully-preserved code, so the trace oracle keeps working; "how far have we
drifted" has a mechanical answer; class extensions are the only tool needed;
and the eventual fork is a directory boundary rather than an archaeology
problem.

Your own code goes in a new directory beside them.

\section{A package}

A package is a directory of Tonel files plus a \ct{package.st}:

\begin{plain}
lib/MyStuff/
    package.st
    Account.class.st
    SavingsAccount.class.st
    Object.extension.st
\end{plain}

\ct{package.st} is one line:

\begin{st}
Package { #name : 'MyStuff' }
\end{st}

\section{A class file}

\ct{Account.class.st}, in full:

\begin{st}
"
The class comment goes at the top of the file, in a Smalltalk comment.
The Browser shows it under the class's `comment' button, so write it for
somebody who has just found this class and does not know why it exists.
"
Class {
	#name : 'Account',
	#superclass : 'Object',
	#instVars : [ 'balance', 'history' ],
	#classVars : [ 'InterestRate' ],
	#category : 'MyStuff-Banking'
}

{ #category : 'initialization' }
Account >> initialize [
	balance := 0.
	history := OrderedCollection new
]

{ #category : 'accessing' }
Account >> balance [
	"How much is in the account."
	^balance
]

{ #category : 'class initialization' }
Account class >> initialize [
	InterestRate := 5 / 100
]
\end{st}

The pieces:

\begin{center}
\small
\begin{tabular}{@{}lp{0.58\textwidth}@{}}
\toprule
\ctp{\#name} \ctp{\#superclass} & required\\
\ctp{\#instVars} \ctp{\#classVars} \ctp{\#classInstVars} & lists of strings\\
\ctp{\#category} & where it shows up in the Browser\\
\ctp{\#traits} \ctp{\#classTraits} & trait composition\\
\ctp{\#type} & \ctp{variable}, \ctp{bytes}, \ctp{words}, \ctp{weak},
  \ctp{ephemeron}. Omit for an ordinary class\\
\bottomrule
\end{tabular}
\end{center}

A method is a category pragma and then the method itself:

\begin{center}
\ct|Class >> selector [ ... ]| \qquad \ct|Class class >> selector [ ... ]|
\end{center}

\noindent the second form being the class side.

\begin{stnote}[Load order does not matter]
Definitions are read in a pass of their own, before any method is compiled. A
class may reference one defined in a later file, or in a later package, and it
loads. This is the one thing the 1983 chunk format could not do, and it is why
there is no dependency declaration anywhere in a package.
\end{stnote}

\section{An extension file}

To add methods to a class you do not own---including 1983's, which you may not
edit---write an extension:

\ct{Object.extension.st}:

\begin{st}
Extension { #name : 'Object' }

{ #category : '*MyStuff' }
Object >> myUsefulThing [
	^42
]
\end{st}

The leading \ct{*} on the category is the convention that marks a method as
belonging to another package. Keep it.

\section{A profile}

A profile says which packages make an image. \ct{profiles/mystuff.profile}:

\begin{st}
"
Free text at the top of the file.  Say why this profile exists and why
anything is excluded -- the profiles in this tree are worth reading as
examples of that.
"
Profile {
	#name     : 'mystuff',
	#requires : [ 'st2026' ],
	#dialect  : 'closures',
	#packages : [ '../lib/MyStuff' ]
}
\end{st}

\begin{center}
\small
\begin{tabular}{@{}lp{0.58\textwidth}@{}}
\toprule
\ctp{\#name} & the profile's name\\
\ctp{\#requires} & other profiles to load first. \ctp{st2026} is the usual
  base; it requires \ctp{bluebook}\\
\ctp{\#dialect} & \ctp{closures} or the Blue Book default\\
\ctp{\#packages} & directories, relative to the profile file\\
\ctp{\#manifests} & a 1983-style manifest, for chunk sources\\
\ctp{\#files} & individual source files\\
\ctp{\#exclude} & class names to drop \emph{wherever} they come from\\
\ctp{\#supersede} & class names this profile provides \emph{instead of} an
  earlier one\\
\bottomrule
\end{tabular}
\end{center}

Build and run it:

\begin{sh}
./st80 -bootstrap -profile profiles/mystuff.profile -eval 'Account new balance'
./st80 -bootstrap -profile profiles/mystuff.profile -o mystuff.image
./st80 -run mystuff.image
\end{sh}

\section{\texttt{\#exclude} versus \texttt{\#supersede}}

The distinction is not pedantic and getting it wrong produces a confusing
image.

\begin{description}
\item[\ctp{\#exclude}] drops a class \emph{wherever} it comes from. The class
  is simply absent. Use it for tooling metadata you do not want, or for a test
  whose subject does not exist here.
\item[\ctp{\#supersede}] says "this profile provides that class instead". The
  earlier definition is dropped and yours is used. Use it when you are
  replacing an implementation.
\end{description}

Excluding something you meant to supersede leaves a hole: every reference to
that class finds \ct{nil}. The \ct{pharo-weak} profile spent a long time at
12 of 32 tests for exactly that reason---\ct{WeakArray} was excluded rather
than superseded, so both providers went and every \ct{WeakSet} had a nil where
its storage should be.

The loader helps: when a supersession drops protocol something still sends, it
reports it, and separates the names nothing else answers---the real holes
---from the far larger number that merely collide with another class's method
of the same name.

\section{Reading the build output}

Every bootstrap prints a report. It is worth understanding rather than
scrolling past.

\begin{plain}
st80: 38 classes were given "new ^super new initialize": ...
st80: 264 classes, 5123 methods, 3840 symbols
st80: 80 undeclared globals, 4 lower-case (probable source bugs): ...
st80: 6 undeclared names resolved from a pool dictionary
st80: 46 class initializers, 43 ran, 3 skipped, 0 unfinished
\end{plain}

\begin{description}
\item[given \ct{new}] classes with an \ct{initialize} and no class-side
  \ct{new}, for which the loader wrote the 1983 idiom. Auditable by design.
\item[undeclared globals] names referenced but never defined. Some are
  expected---the 1983 library has a standing set. A \emph{lower-case} one is
  almost always a typo in source.
\item[class initializers] class-side \ct{initialize} methods. "Unfinished"
  is the one to worry about: an initialiser that started and did not return
  usually means something deadlocked.
\end{description}

\section{Vendoring somebody else's code}

Every imported package carries a \ct{PROVENANCE.md} recording the upstream
repository, the commit, the licence and \emph{every local edit}. That is not
bureaucracy---it is what makes "we import Pharo's collections" a checkable
statement rather than a claim, and it is what lets the next person tell an
upstream bug from one of ours.

If you import something, write one.

\begin{stnote}[Licensing]
The tree is not all one licence. \ct{src/}, \ct{lib/}, \ct{tests/},
\ct{tools/} and \ct{doc/} are BSD 2-Clause. \ct{sources/} is MIT.
\ct{pharo/} is MIT with parts under Apache-2.0, and every file keeps its
notice. The rasterised face derives from Inter under the SIL Open Font
License. \ct{doc/LICENSING.md} is the authority on what may be redistributed.
\end{stnote}
